\chapter{Viewing past data}
\label{ch:viewing}
\index{data viewer}

\section{Opening files}

Press \gui{Load Data File} in the main window, or start the viewer by itself
(Section~\ref{sec:main-window}). Then use \gui{Open File} in the viewer.

\begin{center}
\begin{tabularx}{\textwidth}{@{}lX@{}}
\toprule
File type & What is opened \\
\midrule
\file{.p} & A pyabr3 recording. Opening any one trial loads the whole recording. \\
\file{.txt} & A run from the MATLAB program, ABR4. Any file of the run may be chosen. \\
\file{.pkl} & A threshold summary such as \file{all\_tone\_map\_data.pkl}, or a file of
  numeric arrays. \\
\bottomrule
\end{tabularx}
\end{center}

\begin{caution}
\textbf{One recording per folder.} All of the \file{.p} files in a folder with the same
date and stimulus type are treated as one recording. Keep each recording in its own
folder (Section~\ref{sec:data-files}).
\end{caution}

The viewer is for checking data. Its averages match those of the analysis program, but
it does not measure thresholds or wave amplitudes.

\section{pyabr3 recordings}
\index{recordings!viewing}

The window has settings on the left, the plots on the right, and an \gui{Information}
panel below (Figure~\ref{fig:viewer_clicks}). The \gui{Information} panel shows the
subject data, protocol and calibration stored with the recording.

\figorplaceholder{viewer_clicks}{A stored click recording. The upper plot shows the
average at each level. The middle plot shows one sweep of the raw
data.}{Data viewer with a click recording.}

\begin{center}
\begin{tabularx}{\textwidth}{@{}lX@{}}
\toprule
Setting & Effect \\
\midrule
\gui{High pass}, \gui{Low pass} & Filter corners in Hz; 0 turns a filter off. The
  defaults are 300 and 3000~Hz. \\
\gui{Frequency} & The tone frequency to display \\
\gui{Interleaved set} & For interleaved recordings with more than one stimulus set,
  which set to show \\
\gui{Invert} & Inverts the polarity of the display \\
\gui{Display} & Average only; average with standard deviation; or average with
  individual traces \\
\gui{Stimulus \#}, \gui{Repetition} & Selects the sweep in the \gui{Single trial} plot \\
\bottomrule
\end{tabularx}
\end{center}

The traces are filtered forward and backward with 4-pole Butterworth filters, so no
phase shift is added. The responses are then interpolated to 100~kHz, and converted to
microvolts at the electrodes assuming an amplifier gain of 10,000\index{amplifier gain}.
The first millisecond of each response precedes the stimulus. For click and tone pip
recordings, the responses to all of the stimuli in a sweep are averaged together.

\index{standard deviation}
\gui{Average with SD} shades one standard deviation around the mean
(Figure~\ref{fig:viewer_clicks_sd}). \gui{Average with individual traces} draws up to
100 traces at each level. The first use of either view calculates the single-trial
responses, which takes a short time. The status field shows the progress.

\figorplaceholder{viewer_clicks_sd}{The same recording, with the standard deviation
shown at each level.}{Data viewer showing the average with SD.}

Figure~\ref{fig:viewer_interleaved} shows an interleaved recording. The trace for each
level and frequency is taken from its own part of the sweep.

\figorplaceholder{viewer_interleaved}{An interleaved recording, showing the responses at
16~kHz.}{Data viewer with an interleaved recording.}

The \gui{Stimulus} plot normally says that the waveform is not stored, because later
versions of the program do not save it.

\section{MATLAB (ABR4) runs}
\index{ABR4!viewing runs}\index{MATLAB runs}

A run named \file{yyyymmdd-HHMM} consists of these text files:

\begin{center}
\begin{tabular}{@{}ll@{}}
\toprule
File & Contents \\
\midrule
\file{<run>-SPL.txt} & Levels \\
\file{<run>-kHz.txt} & Frequencies (tone runs only) \\
\file{<run>-p.txt}, \file{<run>-n.txt} & Click responses for each polarity \\
\file{<run>-p-<freq>.txt}, \file{<run>-n-<freq>.txt} & Tone responses for each polarity \\
\bottomrule
\end{tabular}
\end{center}

The \gui{MATLAB run} group lists all of the runs in the folder. The two polarities are
averaged, the first 25~ms are shown, and the filter defaults are 500 and 2500~Hz.
Individual traces are not stored in these files, so only averages can be shown.

\figorplaceholder{viewer_matlab}{A run from the MATLAB program.}{Data viewer with a
MATLAB ABR4 run open. No .txt runs were available in the repository when this draft was
written.}

\begin{caution}
\index{sample rate!MATLAB runs} The sample rate of a MATLAB run is not stored in its
files. \gui{Sample rate (Hz)} is set to 100~kHz, which is correct for recent recordings.
Older recordings may have used 50~kHz. A wrong value changes the time axis and the
effect of the filters.
\end{caution}

\section{Threshold summaries}
\index{thresholds!summary files}

A \file{.pkl} file containing frequencies and thresholds, as written by
\file{read\_abr.py}, is plotted as in Figure~\ref{fig:viewer_thresholds}. Each group
appears as a mean with symbols, and each subject as a faint line. Any other \file{.pkl}
file containing numeric arrays has each array plotted.

\figorplaceholder{viewer_thresholds}{A threshold summary from
\file{all\_tone\_map\_data.pkl}.}{Data viewer showing a threshold summary.}
